% ID: SEL-D0028-OB-E01
% ESTADO: validada. S029, cotejada con examen y criterios.
\textbf{a)} $CA$ tiene dimensión $1\times2$ y $B$ es $2\times1$,
así que $B+2CA$ no existe. En cambio $BC$ es $2\times2$,
al igual que su traspuesta y $A$. Por tanto
\[
A-(BC)^t=
\begin{pmatrix}6&0\\2&4\end{pmatrix}
-\begin{pmatrix}8&-12\\8&-12\end{pmatrix}
=\boxed{\begin{pmatrix}-2&12\\-6&16\end{pmatrix}}.
\]
\textbf{b)} Multiplicamos por 5 y restamos $B$:
$AX=5C^t-B=(-6,-16)^t$. Con $X=(u,v)^t$, las filas dan
$6u=-6$ y $2u+4v=-16$. De ahí
$\boxed{X=\begin{pmatrix}-1\\-7/2\end{pmatrix}}$.
Al sustituir, $B+AX=(-10,-10)^t=5C^t$.
